\section*{Chapitre \ref{ch:b1:complexation} --- Complexation}

\begin{solution}{exo:b1:complexation:1}
$[\ce{Cu(NH3)4}]^{2+}$~: \ce{Cu^2+}, quatre \ce{NH3} liés par N.
$[\ce{Fe(CN)6}]^{4-}$~: \ce{Fe^2+}, six \ce{CN-} liés par C~;
$[\ce{Fe(CN)6}]^{3-}$~: de même avec \ce{Fe^3+}. $[\ce{Ag(NH3)2}]^{+}$~:
\ce{Ag+}, deux \ce{NH3} liés par N. $[\ce{FeSCN}]^{2+}$~: \ce{Fe^3+}, un
\ce{SCN-}, lié par N malgré la façon dont la formule est écrite.
$[\ce{Zn(OH)4}]^{2-}$~: \ce{Zn^2+}, quatre \ce{OH-} liés par O.
$[\ce{CaY}]^{2-}$~: \ce{Ca^2+}, un ion edta lié par deux N et quatre O.
\end{solution}

\begin{solution}{exo:b1:complexation:2}
$\log K_i = \log\beta_i - \log\beta_{i-1}$~: 4.10, 3.41, 2.82, 2.03, et
$\mathrm{p}K_{d,i} = \log K_i$ prend les mêmes valeurs. Le dernier se
rapporte à
$[\ce{Cu(NH3)4}]^{2+} \rightleftharpoons [\ce{Cu(NH3)3}]^{2+} + \ce{NH3}$.
\end{solution}

\begin{solution}{exo:b1:complexation:3}
Domaines~: $[\ce{Cu(NH3)4}]^{2+}$ en dessous de pNH3 2.03, puis $n = 3$
jusqu'à 2.82, $n = 2$ jusqu'à 3.41, $n = 1$ jusqu'à 4.10, \ce{Cu^2+}
au-dessus. pNH3 = 0~:
$[\ce{Cu(NH3)4}]^{2+}$~; 2.5~: $[\ce{Cu(NH3)3}]^{2+}$~; 3.5~:
$[\ce{CuNH3}]^{2+}$~; 5.0~: \ce{Cu^2+}.
\end{solution}

\begin{solution}{exo:b1:complexation:4}
$[\ce{FeSCN^2+}]/[\ce{Fe^3+}] = \beta_1[\ce{SCN-}] = 10^{2.96} \times 0.10 =
91$~: la fraction complexée vaut $91/92 = \qty{98.9}{\percent}$.
\end{solution}

\begin{solution}{exo:b1:complexation:5}
Chaque $[\mathrm{ML}_k] = \beta_k[\mathrm{M}][\mathrm{L}]^k$~; en divisant
par la somme sur $k$ (le métal total), on élimine $[\mathrm{M}]$. Avec
trois espèces, $f_i = 1/(1 + [\mathrm{ML}_{i-1}]/[\mathrm{ML}_i] +
[\mathrm{ML}_{i+1}]/[\mathrm{ML}_i]) = 1/(1 + 1/(K_i x) + K_{i+1}x)$ avec
$x = [\mathrm{L}]$. Le dénominateur est minimal quand sa dérivée
$-1/(K_i x^2) + K_{i+1}$ s'annule, $x^2 = 1/(K_iK_{i+1})$, là où les deux
voisins sont égaux, et $\mathrm{pL} = \frac12(\log K_i + \log K_{i+1})$.
Pour $[\ce{Cu(NH3)2}]^{2+}$~: $\frac12(3.41 + 2.82) = 3.12$.
\end{solution}

\begin{solution}{exo:b1:complexation:6}
$[\ce{FeSCN}]^{2+} + \ce{F-} \rightleftharpoons [\ce{FeF}]^{2+} + \ce{SCN-}$,
$K = 10^{6.85 - 2.96} = 10^{3.89}$. Rapport $= K[\ce{F-}]/[\ce{SCN-}] =
10^{3.89} \times 0.010/0.10 = 780$~: le \omterm{def:b1:complexation:complex}{complexe} rouge est presque
entièrement remplacé et la couleur s'efface.
\end{solution}

\begin{solution}{exo:b1:complexation:7}
$[\ce{MgY}]^{2-} + \ce{Ca^2+} \rightleftharpoons [\ce{CaY}]^{2-} + \ce{Mg^2+}$,
$K = 10^{12.69 - 10.90} = 10^{1.79} = 62$. Avec des quantités égales,
$x^2/(1 - x)^2 = 62$, $x/(1 - x) = 7.9$, $x = 0.89$~: \qty{89}{\percent}
du magnésium est libéré.
\end{solution}

\begin{solution}{exo:b1:complexation:8}
\ce{Ag2O(s) + H2O + 4NH3 <=> 2[Ag(NH3)2]+ + 2OH-}, $K = \beta_2^2 \times
10^{-15.43} = 10^{14.44 - 15.43} = 10^{-0.99}$. La constante n'est pas
petite~: dès que l'ammoniac libre atteint quelques dixièmes de mole par
litre, l'oxyde brun formé par les premières gouttes (l'ammoniac agissant
alors comme une base) se dissout en donnant le \omterm{def:b1:complexation:complex}{complexe} ammine de
l'argent.
\end{solution}

\begin{solution}{exo:b1:complexation:9}
pH 10~: $\alpha = 1 + 10^{1.24} + 10^{-1.96} = 18.4$, $\log\beta' = 12.69 -
1.26 = 11.43$. pH 7~: $\alpha = 1 + 10^{4.24} + 10^{4.04} + 10^{0.19} +
\cdots = 10^{4.45}$, $\log\beta' = 8.24$. À pH 7, la constante
conditionnelle est inférieure à $10^{10}$~: la réaction du calcium avec
l'edta n'est pas assez complète pour un titrage, parce que l'edta est en
majeure partie protoné.
\end{solution}

\begin{solution}{exo:b1:complexation:10}
1.~\ce{CuO(s) + H2O + 4NH3 <=> [Cu(NH3)4]^2+ + 2OH-}, $K = \beta_4 \times
10^{-20.64} = 10^{12.36 - 20.64} = 10^{-8.28}$.
2.~$\mathrm{pH} = 7 + \frac12(9.25 + \log 1.0) = 11.63$, $[\ce{OH-}] =
10^{-2.37}$, et $[\ce{Cu(NH3)4^2+}] = 10^{-8.28}/10^{-4.75} =
\qty{3.0e-4}{mol/L}$~: très peu de solide se dissout.
3.~$\mathrm{pH} = 9.25$, $[\ce{OH-}] = 10^{-4.75}$, et la formule donne
$10^{-8.28 + 9.50} = 10^{1.22}$, bien plus que ce que l'ammoniac peut
complexer~: l'oxyde se dissout jusqu'à épuisement de l'ammoniac, et non
jusqu'à l'équilibre. Les ions ammonium consomment les ions hydroxyde
libérés par la dissolution~; c'est pourquoi les sels d'ammonium aident
l'ammoniac à dissoudre les composés du cuivre.
\end{solution}

\begin{solution}{exo:b1:complexation:11}
1.~La réaction est la deuxième formation moins la première~:
$K = K_2/K_1 =
10^{3.90 - 3.32} = 10^{0.58} = 3.8 > 1$, si bien que $[\ce{AgNH3}]^{+}$
se dismute lorsqu'il est l'espèce majoritaire.
2.~D'après l'\cref{exo:b1:complexation:5}, la plus grande fraction est
atteinte en
$\mathrm{pNH_3} = \frac12(3.32 + 3.90) = 3.61$, où $\beta_1 x = 10^{-0.29}
= 0.51$ et $\beta_2x^2 = 1.0$~: $f = 0.51/(1 + 0.51 + 1.0) = 0.20$. Jamais
plus de \qty{20}{\percent} de l'argent ne se trouve sous forme de
$[\ce{AgNH3}]^{+}$.
\end{solution}

\begin{solution}{exo:b1:complexation:12}
1.~\ce{CaCO3(s) + Y^4- <=> [CaY]^2- + CO3^2-}, $K = \beta K_s = 10^{12.69 -
8.30} = 10^{4.39}$.
2.~$x^2/(0.10 - x) = 10^{4.39}$ donne $x = \qty{0.10}{mol/L}$ à
$4 \times 10^{-7}$ près~: tout l'edta est utilisé, soit \qty{10.0}{g} de
carbonate de calcium par litre.
3.~En milieu acide, l'edta est protoné ($\mathrm{p}K_{a4} = 11.24$) et
sa constante conditionnelle chute~; le carbonate est lui aussi protoné,
ce qui aide, mais c'est le \omterm{def:b1:complexation:complex}{complexe} qui retient le calcium. Maintenir un
pH élevé garde l'edta sous la forme \ce{Y^4-}.
\end{solution}

\begin{solution}{pb:b1:complexation:1}
\textbf{1.} \ce{Ag+ + NH3 <=> [AgNH3]+}, $\beta_1 = [\ce{AgNH3+}]/([\ce{Ag+}][\ce{NH3}])$;
\ce{Ag+ + 2NH3 <=> [Ag(NH3)2]+}, $\beta_2 = [\ce{Ag(NH3)2+}]/([\ce{Ag+}][\ce{NH3}]^2)$.
\textbf{2.} $\log K_1 = 3.32$, $\log K_2 = 7.22 - 3.32 = 3.90$.
\textbf{3.} $[\ce{AgNH3}]^{+}$ prédominerait entre pNH3 3.90 et 3.32,
un intervalle à l'envers~: il n'a pas de domaine.
\textbf{4.} \ce{2[AgNH3]+ <=> Ag+ + [Ag(NH3)2]+}, $K = K_2/K_1 = 10^{0.58} =
3.8$.
\textbf{5.} $[\ce{Ag(NH3)2}]^{+}$ en dessous et \ce{Ag+} au-dessus de
$\frac12\log\beta_2 = 3.61$.
\textbf{6.} $\beta_2[\ce{NH3}]^2 = 10^{7.22 - 2.00} = 10^{5.22} = \num{1.7e5}$.
\textbf{7.} $\mathrm{p}K_d = \log\beta_2 = 7.22$.
\textbf{8.} \ce{AgCl(s) + 2NH3 <=> [Ag(NH3)2]+ + Cl-}, $K = \beta_2K_s =
10^{7.22 - 9.75} = 10^{-2.53}$.
\textbf{9.} $s = \sqrt{K}\,[\ce{NH3}] = 10^{-1.265} \times 1.0 =
\qty{0.054}{mol/L}$.
\textbf{10.} Dans l'eau, $s = 10^{-4.875} = \qty{1.3e-5}{mol/L}$~:
environ 4000 fois plus.
\textbf{11.} $0.0543 \times 143.4 = \qty{7.8}{g}$.
\textbf{12.} $[\ce{Ag+}] = K_s/[\ce{Cl-}] = 10^{-9.75}/0.054 =
\qty{3.3e-9}{mol/L} \ll s$.
\textbf{13.} $s = \sqrt{K}(1.0 - 2s)$, $s = 0.0543/(1 + 2 \times 0.0543) =
\qty{0.049}{mol/L}$.
\textbf{14.} L'ammoniac à \qty{1.0}{mol/L} a un pH de 11.6, supérieur à
$9.25 + 1$~: seulement \qty{0.4}{\percent} est sous forme \ce{NH4+}.
\textbf{15.} $K = 10^{7.22 - 12.27} = 10^{-5.05}$.
\textbf{16.} $s = 10^{-2.525} = \qty{3.0e-3}{mol/L}$, \qty{0.56}{g/L}.
\textbf{17.} $K = 10^{-8.85}$, $s = 10^{-4.425} = \qty{3.8e-5}{mol/L}$,
\qty{8.8}{mg/L}.
\textbf{18.} Traiter chaque précipité par de l'ammoniac dilué~: seul le
chlorure se dissout. L'ammoniac concentré dissout le bromure, plus
lentement et en partie~; l'iodure ne se dissout pas.
\textbf{19.} $s = 1.0/187.8 = \qty{5.3e-3}{mol/L}$~; $[\ce{NH3}] = s/\sqrt{K}
= \num{5.3e-3}/10^{-2.525} = \qty{1.8}{mol/L}$.
\textbf{20.} $s = \qty{4.3e-3}{mol/L}$ et $[\ce{NH3}] = \num{4.3e-3}/10^{-4.425}
= \qty{113}{mol/L}$, deux fois la concentration de l'eau dans l'eau pure
(\qty{55}{mol/L})~: impossible. L'ammoniac ne peut pas dissoudre
l'iodure d'argent.
\textbf{21.} \ce{[Ag(NH3)2]+ + Cl- + 2H3O+ -> AgCl(s) + 2NH4+ + 2H2O},
$K = 10^{2.53} \times (10^{9.25})^2 = 10^{21.03}$~: le précipité blanc
réapparaît, ce qui confirme le chlorure d'argent.
\textbf{22.} $s = 1.0/143.4 = \qty{6.97e-3}{mol/L}$~; $[\ce{NH3}] =
s/\sqrt{K} = \num{6.97e-3}/10^{-1.265} = \qty{0.128}{mol/L}$.
\textbf{23.} $2s = \qty{0.0139}{mol/L}$.
\textbf{24.} $[\ce{Ag+}] = 10^{-9.75}/\num{6.97e-3} = \qty{2.6e-8}{mol/L}$,
négligeable devant $s$.
\textbf{25.} $0.128 + 0.014 = $ \textbf{\qty{0.142}{mol/L} d'ammoniac}~:
\qty{0.142}{mol} dans le litre dissout exactement \qty{1.0}{g} de
chlorure d'argent.
\end{solution}
